\documentclass[11pt]{article}

\usepackage[margin=1in]{geometry}
\usepackage{amsmath, amssymb}
\usepackage{xcolor}
\usepackage{listings}
\usepackage{hyperref}
\usepackage{array}
\usepackage{booktabs}

\hypersetup{
    colorlinks=true,
    linkcolor=blue!60!black,
    urlcolor=blue!60!black,
    citecolor=blue!60!black
}

\lstdefinestyle{bugpython}{
  language=Python,
  basicstyle=\ttfamily\small,
  keywordstyle=\color{blue!60!black},
  stringstyle=\color{green!40!black},
  commentstyle=\color{gray!70!black},
  showstringspaces=false,
  breaklines=true,
  frame=single,
  rulecolor=\color{gray!30},
  columns=fullflexible,
  keepspaces=true
}

\title{SymPy Bug Report}
\author{}
\date{}

\begin{document}

\maketitle

\section{Bug 12: Inverse Fourier transform drops Abs in the Lorentzian pair}

\subsection{Minimal Reproducer}

\medskip

\begin{lstlisting}[style=bugpython]
from sympy import exp, inverse_fourier_transform, pi, simplify, symbols

t, w = symbols("t w", real=True)
F = 2/(1 + 4*pi**2*w**2)
inv = inverse_fourier_transform(F, w, t)

print("inverse =", inv)
print("at t=-1:", inv.subs(t, -1))
print("expected at t=-1:", exp(-1))
print("difference:", simplify(inv.subs(t, -1) - exp(-1)))
\end{lstlisting}

\subsection{Observed Behavior}

\medskip

\begin{lstlisting}[style=bugpython]
SymPy version: 1.14.0
inverse = exp(-t)
at t=-1: E
expected at t=-1: exp(-1)
difference: 2*sinh(1)
\end{lstlisting}

SymPy returns \(e^{-t}\) as an unconditional inverse Fourier transform.  This
expression agrees with the correct inverse only on the nonnegative real half
line.  At \(t=-1\), the returned expression evaluates to \(e\), while the
expected value is \(e^{-1}\).

\subsection{Expected Behavior}

The inverse Fourier transform should be
\[
  \operatorname{inverse\_fourier\_transform}
  \left(\frac{2}{1+4\pi^2 w^2}, w, t\right)
  = e^{-|t|}.
\]
Equivalently, for real \(t\), the result should be the even function
\(\exp(-\operatorname{Abs}(t))\), not the one-sided branch \(\exp(-t)\).

\subsection{Mathematical Explanation}

With SymPy's Fourier convention,
\[
  \mathcal{F}\{f\}(w)
  = \int_{-\infty}^{\infty} f(t)e^{-2\pi iwt}\,dt,
  \qquad
  \mathcal{F}^{-1}\{F\}(t)
  = \int_{-\infty}^{\infty} F(w)e^{2\pi iwt}\,dw .
\]
The standard Lorentzian transform pair is
\[
  \int_{-\infty}^{\infty} e^{-a|t|}e^{-2\pi iwt}\,dt
  = \frac{2a}{a^2+4\pi^2w^2}
  \qquad (a>0).
\]
For \(a=1\), the inverse transform of
\(2/(1+4\pi^2w^2)\) is therefore \(e^{-|t|}\).  The expression \(e^{-t}\) is
not an equivalent representation of \(e^{-|t|}\): for \(t<0\),
\[
  e^{-t}=e^{|t|}\neq e^{-|t|}.
\]
The discrepancy is not a missing multiplicative constant or a harmless branch
choice.  It changes the decay on the negative half line into exponential growth.

\subsection{Source-Level Diagnosis}

The diagnosis identifies the relevant implementation as
\nolinkurl{sympy/integrals/transforms.py}, specifically
\texttt{\_fourier\_transform} around lines 945--961.  The public call
\texttt{inverse\_fourier\_transform(F, w, t)} constructs an inverse transform
object.  Its compute method follows the shared Fourier-transform path and then
calls \texttt{\_fourier\_transform} with the inverse-Fourier parameters.
That helper evaluates the real-axis kernel integral
\[
  \int_{-\infty}^{\infty} F(w)e^{2\pi iwt}\,dw .
\]

For the reproducer, the raw integral result is a \texttt{Piecewise} value whose
first branch is a closed form valid only under the condition \(t>0\), with an
unevaluated integral as the fallback branch.  The diagnosed problematic rule is
that \texttt{\_fourier\_transform} accepts a \texttt{Piecewise} integral result
by taking only the first branch and returning that branch expression together
with its condition:

\begin{lstlisting}[style=bugpython]
if F.is_Piecewise:
    F, cond = F.args[0]
    return _simplify(F), cond
\end{lstlisting}

In this case, the first branch simplifies to \(\exp(-t)\) and its condition is
\(t>0\).  The public Fourier transform wrapper uses \texttt{noconds=True} by
default, so the condition is discarded before the result is returned to the
caller.  The source-level failure is therefore a conditional branch being
promoted to an unconditional transform result.  A plausible fix direction,
supported by the diagnosis, is to avoid collapsing a conditional
\texttt{Piecewise} transform integral to its first branch unless that condition
is globally valid; otherwise SymPy should preserve the conditional result,
compute the missing branch, or leave the transform unevaluated when
\texttt{noconds=True} would hide a necessary condition.

\subsection{Additional Failing Instances}

The same error appears for the parametrized Lorentzian family
\[
  F_a(w)=\frac{2a}{a^2+4\pi^2w^2}, \qquad a>0.
\]
Mathematically, this is the Fourier transform of \(e^{-a|t|}\), so every
inverse transform in the family should be even in \(t\).  The verification file
checks \(a=1,2,3,4,5,6\) at \(t=-1\); in each case the returned positive-\(t\)
branch gives \(e^a\) rather than \(e^{-a}\).

\begin{center}
\small
\renewcommand{\arraystretch}{1.3}
\begin{tabular}{@{}>{\raggedright\arraybackslash}p{0.44\linewidth}>{\raggedright\arraybackslash}p{0.23\linewidth}>{\raggedright\arraybackslash}p{0.23\linewidth}@{}}
\toprule
\textbf{Input / Expression} & \textbf{SymPy behavior} & \textbf{Expected behavior} \\
\midrule
\(2/(1+4\pi^2w^2)\) & returns \(e^{-t}\); at \(t=-1\), \(e\) & \(e^{-|t|}\); at \(t=-1\), \(e^{-1}\) \\
\(4/(4+4\pi^2w^2)\) & returns \(e^{-2t}\); at \(t=-1\), \(e^2\) & \(e^{-2|t|}\); at \(t=-1\), \(e^{-2}\) \\
\(6/(9+4\pi^2w^2)\) & returns \(e^{-3t}\); at \(t=-1\), \(e^3\) & \(e^{-3|t|}\); at \(t=-1\), \(e^{-3}\) \\
\(8/(16+4\pi^2w^2)\) & returns \(e^{-4t}\); at \(t=-1\), \(e^4\) & \(e^{-4|t|}\); at \(t=-1\), \(e^{-4}\) \\
\(10/(25+4\pi^2w^2)\) & returns \(e^{-5t}\); at \(t=-1\), \(e^5\) & \(e^{-5|t|}\); at \(t=-1\), \(e^{-5}\) \\
\(12/(36+4\pi^2w^2)\) & returns \(e^{-6t}\); at \(t=-1\), \(e^6\) & \(e^{-6|t|}\); at \(t=-1\), \(e^{-6}\) \\
\bottomrule
\end{tabular}
\end{center}

\subsection{Related Issues Assessment}

The bug-finding harness performed a best-effort deduplication check and found
no public issue, pull request, release-note entry, Stack Overflow post, or
mailing-list thread describing this \emph{specific} Lorentzian inverse transform
returning \(\exp(-t)\) instead of \(\exp(-\operatorname{Abs}(t))\).  The
underlying failure family, however, is already documented in the same
\texttt{\_fourier\_transform} code path.  SymPy
\href{https://github.com/sympy/sympy/issues/22787}{issue \#22787}
(\emph{Bug in fourier\_transform}, open, still reproducing on 1.14.0) reports
\texttt{fourier\_transform}/\texttt{\_fourier\_transform} returning a non-even
result for a real even input, using the same evenness argument, and
\href{https://github.com/sympy/sympy/issues/14624}{issue \#14624} (closed)
reports \texttt{inverse\_fourier\_transform} giving a wrong answer on the
negative half-line.  The deduplication verdict is therefore
\emph{family known, specific instance new}: this is a new concrete instance of
the known, still-open \texttt{\_fourier\_transform} Piecewise-branch failure
family, not a novel bug.  The case remains individually actionable because it
gives a concrete, reproducible transform-pair counterexample and a narrow
regression test, and the regression test may be linked to issue \#22787.

\subsection{Proposed SymPy Regression Test}

\medskip

\begin{lstlisting}[style=bugpython]
import pytest
from sympy import Abs, exp, inverse_fourier_transform, pi, simplify, symbols


@pytest.mark.parametrize("a", [1, 2, 3, 4, 5, 6])
def test_inverse_fourier_lorentzian_is_even(a):
    t, w = symbols("t w", real=True)
    expr = 2*a/(a**2 + 4*pi**2*w**2)
    inv = inverse_fourier_transform(expr, w, t)
    assert simplify(inv.subs(t, -1) - exp(-a*Abs(-1))) == 0
\end{lstlisting}

\end{document}
